\documentclass[11pt]{article}
\usepackage[utf8]{inputenc}
\usepackage[T1]{fontenc}
\usepackage[spanish]{babel}
\usepackage{amsmath}
\usepackage{booktabs}
\usepackage{graphicx}
\usepackage{float}
\usepackage{geometry}
\geometry{margin=2.5cm}
\setlength{\parskip}{0.65em}
\setlength{\parindent}{0pt}
\title{Testing probabilidad de abandono: Version 3}
\author{FinPUC}
\date{}
\begin{document}
\maketitle
\section{Análisis de resultados fuera de muestra}
La validación financiera se replica desde la iteración de desempleo macro, manteniendo fija la arquitectura de tres horizontes y la calibración Black-Litterman. El objetivo de esta carpeta no es volver a seleccionar una view, sino aislar el efecto de la función de abandono sobre la capa económica P4.
\[
P_{retiro,t}=
\begin{cases}
\dfrac{0.10}{1+\exp[-20(loss_t-\hat{x}_1)]}, & loss_t>\hat{x}_1,\\
0, & loss_t\leq \hat{x}_1.
\end{cases}
\]
Esta versión combina techo de 10% semanal con una pendiente mayor. Con ello, el abandono se mantiene acotado, pero reacciona más rápido cuando la pérdida acumulada excede con claridad la tolerancia declarada del perfil.
\begin{table}[H]
\centering
\resizebox{\textwidth}{!}{%
\begin{tabular}{lrrrrrrr}
\toprule
Escenario & Perfil & Sharpe BL & Sharpe MK & Mejora Sharpe & DD BL & DD MK & \% recomendado \\
\midrule
sin\_pandemia & Muy conservador & 1.738 & 1.736 & 0.2\% & -10.2\% & -10.3\% & 100.0\% \\
sin\_pandemia & Conservador & 1.718 & 1.728 & -0.6\% & -10.1\% & -11.0\% & 0.0\% \\
sin\_pandemia & Neutro & 1.664 & 1.747 & -4.8\% & -10.2\% & -13.0\% & 0.0\% \\
sin\_pandemia & Arriesgado & 1.562 & 1.513 & 3.2\% & -12.1\% & -18.2\% & 100.0\% \\
sin\_pandemia & Muy arriesgado & 1.154 & 1.198 & -3.7\% & -21.1\% & -26.5\% & 0.0\% \\
con\_pandemia & Muy conservador & 1.625 & 1.699 & -4.4\% & -9.1\% & -9.0\% & 0.0\% \\
con\_pandemia & Conservador & 1.638 & 1.749 & -6.5\% & -9.2\% & -9.0\% & 0.0\% \\
con\_pandemia & Neutro & 1.660 & 1.773 & -6.0\% & -9.4\% & -10.9\% & 25.0\% \\
con\_pandemia & Arriesgado & 1.644 & 1.391 & 19.9\% & -11.1\% & -18.7\% & 75.0\% \\
con\_pandemia & Muy arriesgado & 1.119 & 0.933 & 23.1\% & -22.8\% & -27.6\% & 100.0\% \\
\bottomrule
\end{tabular}
}
\caption{Comparación financiera fuera de muestra para Version 3.}
\label{tab:version_3:oos}
\end{table}

\section{Resultados de la dinámica de portafolio}
La dinámica de portafolio se mantiene como control metodológico: cualquier diferencia material entre versiones debe provenir de la simulación de abandono y no de cambios en pesos, rebalanceos o concentración sectorial.
\begin{table}[H]
\centering
\resizebox{\textwidth}{!}{%
\begin{tabular}{lrrrrr}
\toprule
Escenario & Perfil & Turnover & N efectivo & HHI sector & Peso sector líder \\
\midrule
sin\_pandemia & Muy conservador & 7.0\% & 221.21 & 0.147 & 27.5\% \\
sin\_pandemia & Conservador & 7.0\% & 233.51 & 0.144 & 27.2\% \\
sin\_pandemia & Neutro & 6.9\% & 271.13 & 0.137 & 26.2\% \\
sin\_pandemia & Arriesgado & 6.9\% & 360.45 & 0.124 & 21.8\% \\
sin\_pandemia & Muy arriesgado & 9.2\% & 163.27 & 0.178 & 32.1\% \\
con\_pandemia & Muy conservador & 6.3\% & 112.55 & 0.169 & 25.7\% \\
con\_pandemia & Conservador & 6.2\% & 118.08 & 0.166 & 25.0\% \\
con\_pandemia & Neutro & 6.4\% & 136.34 & 0.158 & 23.5\% \\
con\_pandemia & Arriesgado & 7.0\% & 205.71 & 0.135 & 20.1\% \\
con\_pandemia & Muy arriesgado & 9.9\% & 120.96 & 0.214 & 37.9\% \\
\bottomrule
\end{tabular}
}
\caption{Dinámica de portafolio de la view desempleo macro bajo Version 3.}
\label{tab:version_3:dinamica}
\end{table}

\section{Evaluación Económica Monte Carlo P4 en Horizonte Limpio}
La simulación P4 usa 2000 trayectorias por cruce modelo-escenario-perfil y un horizonte de 260 semanas. La utilidad de la empresa se calcula ex-post con una comisión de 0,5\% mensual sobre el saldo administrado activo. La probabilidad de abandono se evalúa semanalmente, pero se reporta también como probabilidad semestral para la presentación.
En el promedio de escenarios y perfiles, la probabilidad semestral de abandono es 0.9\%, la tasa de retiro realizada en 260 semanas es 4.0\%, la retención de clientes activos es 92.7\% y la utilidad media de la empresa alcanza 245. Estos resultados permiten evaluar si Version 3 corrige la saturación observada en la iteración previa sin producir una permanencia artificial de clientes ante pérdidas materialmente superiores a su tolerancia.
\begin{table}[H]
\centering
\resizebox{\textwidth}{!}{%
\begin{tabular}{lrrrrrrrrr}
\toprule
Escenario & Perfil & Riqueza & Utilidad & P acepta reb. & P abandono semanal & P abandono semestre & Retiro 260w & Retención & Clientes finales \\
\midrule
sin\_pandemia & Muy conservador & 1.765 & 222 & 55.7\% & 0.1\% & 3.7\% & 15.8\% & 71.8\% & 805 \\
sin\_pandemia & Conservador & 1.887 & 263 & 54.3\% & 0.0\% & 0.8\% & 4.0\% & 92.7\% & 997 \\
sin\_pandemia & Neutro & 1.864 & 265 & 51.6\% & 0.0\% & 0.0\% & 0.1\% & 99.7\% & 1030 \\
sin\_pandemia & Arriesgado & 1.791 & 244 & 47.9\% & 0.0\% & 0.0\% & 0.0\% & 100.0\% & 957 \\
sin\_pandemia & Muy arriesgado & 1.876 & 252 & 46.0\% & 0.0\% & 0.0\% & 0.2\% & 99.6\% & 943 \\
con\_pandemia & Muy conservador & 1.602 & 198 & 54.9\% & 0.2\% & 4.0\% & 16.3\% & 70.4\% & 776 \\
con\_pandemia & Conservador & 1.752 & 248 & 53.7\% & 0.0\% & 0.7\% & 3.6\% & 93.2\% & 993 \\
con\_pandemia & Neutro & 1.780 & 252 & 51.3\% & 0.0\% & 0.0\% & 0.1\% & 99.7\% & 1020 \\
con\_pandemia & Arriesgado & 1.807 & 247 & 47.9\% & 0.0\% & 0.0\% & 0.0\% & 100.0\% & 972 \\
con\_pandemia & Muy arriesgado & 1.939 & 255 & 46.2\% & 0.0\% & 0.0\% & 0.2\% & 99.6\% & 916 \\
\bottomrule
\end{tabular}
}
\caption{Resultados P4 y comportamiento de clientes para Version 3.}
\label{tab:version_3:p4}
\end{table}

\begin{figure}[H]
\centering
\includegraphics[width=0.88\textwidth]{figuras/fig_probabilidades_sin_pandemia.png}
\caption{Probabilidades de aceptación y abandono semestral sin pandemia para Version 3.}
\label{fig:version_3:prob_sin}
\end{figure}
\begin{figure}[H]
\centering
\includegraphics[width=0.88\textwidth]{figuras/fig_probabilidades_con_pandemia.png}
\caption{Probabilidades de aceptación y abandono semestral con pandemia para Version 3.}
\label{fig:version_3:prob_con}
\end{figure}
\section{Discusión de Robustez}
La robustez se interpreta como estabilidad económica y conductual bajo escenarios con y sin pandemia. Una función de abandono demasiado agresiva reduce artificialmente el saldo administrado y deprime la utilidad; una función demasiado laxa puede esconder pérdidas relevantes al mantener clientes que deberían retirarse según su tolerancia.
\section{Conclusiones}
\begin{table}[H]
\centering
\resizebox{\textwidth}{!}{%
\begin{tabular}{lrrrrrr}
\toprule
Perfil & P abandono semestre & Retiro 260w & Retención & Clientes finales & Utilidad & Lectura \\
\midrule
Muy conservador & 3.9\% & 16.0\% & 71.1\% & 791 & 210 & Estable \\
Conservador & 0.8\% & 3.8\% & 92.9\% & 995 & 255 & Estable \\
Neutro & 0.0\% & 0.1\% & 99.7\% & 1025 & 258 & Estable \\
Arriesgado & 0.0\% & 0.0\% & 100.0\% & 965 & 246 & Estable \\
Muy arriesgado & 0.0\% & 0.2\% & 99.6\% & 930 & 254 & Estable \\
\bottomrule
\end{tabular}
}
\caption{Conclusión segmentada por perfil para Version 3.}
\label{tab:version_3:conclusiones}
\end{table}

La versión debe seleccionarse según su capacidad de evitar la saturación semestral y, al mismo tiempo, mantener sensibilidad ante pérdidas por encima de la tolerancia. En particular, la comparación entre versiones permite defender una especificación conductual calibrada por estabilidad de clientes y no solo por riqueza terminal.
\end{document}